%# -*- coding: utf-8-unix -*-
%%==================================================
%% thesis.tex
%%==================================================

% 双面打印
\documentclass[bachelor, openright, twoside]{zstuthesis}
% \documentclass[bachelor, openany, oneside, submit]{zstuthesis}
% \documentclass[master, review]{zstuthesis}
% \documentclass[%
%   bachelor|master|doctor|coursepaper, % 必选项，分别是学士，硕士，博士学位论文以及课程论文
%   fontset=fandol|windows|mac|ubuntu|adobe|founder, % 字体选项
%   oneside|twoside,        % 单面打印，双面打印(奇偶页交换页边距，默认)
%   openany|openright,      % 可以在奇数或者偶数页开新章|只在奇数页开新章(默认)
%   english,                % 启用英文模版
%   review,     % 盲审论文，隐去作者姓名、学号、导师姓名、致谢、发表论文和参与的项目
%   submit      % 定稿提交的论文，插入签名扫描版的原创性声明、授权声明
% ]

% 逐个导入参考文献数据库
\addbibresource{bib/thesis.bib}
% \addbibresource{bib/chap2.bib}

%# -*- coding: utf-8-unix -*-
% !TEX program = xelatex
% !TEX root = ../thesis.tex
% !TEX encoding = UTF-8 Unicode
%TC:ignore
\title{浙江理工大学学位论文 \\ \LaTeX{} 模板示例文档}
\author{某\quad{}某}
\advisor{某某教授}
% \coadvisor{某某教授}
\defenddate{2019年4月13日}
\coursename{某某课程}
\school{浙江理工大学}
\institute{某某系}
\studentnumber{0010900990}
\cnacademicdegree{工学硕士}
\major{某某专业}
\keywords{浙江理工, 五个浙理, 爱国荣校}

\englishtitle{A Sample Document for \LaTeX-based ZSTU Thesis Template}
\englishauthor{\textsc{Mo Mo}}
\englishadvisor{Prof. \textsc{Mou Mou}}
% \englishcoadvisor{Prof. \textsc{Uom Uom}}
\englishschool{Zhejiang Sci-Tech University}
\englishinstitute{\textsc{Depart of XXX, School of XXX} \\
  \textsc{Zhejiang Sci-Tech University} \\
  \textsc{Shanghai, P.R.China}}
\englishinstitutemaster{Depart of XXX, \\ School of XXX}
\englishmajor{A Very Important Major}
\englishdate{Apr. 13th, 2019}
\enacademicdegree{Master of Engineering}
\englishstudentid{0010900990}
\englishkeywords{ZSTU, master thesis, XeTeX/LaTeX template}
%TC:endignore
  % NOTE: the enclosed commands must be executed in preamble

\begin{document}

% 无编号内容：中英文论文封面、授权页
\maketitle
\makeatletter

\ifzstu@coursepaper
% 摘要 部分课程论文需要，可以自行选择添加或者去除
% %# -*- coding: utf-8-unix -*-
% !TEX program = xelatex
% !TEX root = ../thesis.tex
% !TEX encoding = UTF-8 Unicode
%%==================================================
%% abstract.tex for SJTU Master Thesis
%%==================================================

\begin{abstract}

学校办学历史悠久。浙江理工大学坐落在历史文化名城杭州市，占地面积2100余亩。学校前身——蚕学馆，是杭州知府林启为实现实业救国、教育救国的宏愿于1897年创办的，是我国最早创办的新学教育机构之一。1908年，因办学成绩卓著，被清政府御批升格为“高等蚕桑学堂”。辛亥革命至解放前夕，因时局动乱，学校几度易名、数迁校址，在风雨沧桑中生生不息、坚持办学。新中国成立后，学校不断开拓进取，绘就了新的历史篇章。学校1959年开始招收本科生，1964年由国务院定名为浙江丝绸工学院，1979年开始招收硕士研究生，1983年获硕士学位授予权。1999年，经教育部批准，学校更名为浙江工程学院。2004年，经教育部批准，学校更名为浙江理工大学。2006年，获得博士学位授予权。2010年，经教育部批准，获得推荐优秀应届本科毕业生免试攻读硕士学位研究生资格。2017年，学校入选浙江省重点建设高校。

今天的浙江理工大学是一所以工为主，特色鲜明，优势突出，理、工、文、经、管、法、艺术、教育等多学科协调发展的省重点建设高校。学校下设17个学院（教研部），举办1所独立学院，现有全日制在校学生27000余人，其中研究生3800余人。学校现有本科专业65个，其中国家（教育部）特色专业8个，教育部综合改革试点专业1个、教育部“卓越工程师教育培养计划”试点专业6个、“教育部-欧特克公司专业综合改革项目”3个，“十三五”省优势专业7个、省特色专业5个。拥有省重点建设高校优势特色学科2个，博士学位授权一级学科3个、硕士学位授权一级学科26个，并可授予专业学位类别10种，博士后科研流动站2个。拥有省高校重中之重一级学科1个、省高校重中之重（一级）学科3个、省高校人文社会科学重点研究基地1个、省高校重点学科11个，省一流学科A类6个、一流学科B类7个，材料科学、化学、工程学3个学科进入ESI全球排名前1\%；拥有2个国家地方联合工程实验室，1个国家地方联合工程研究中心，3个国家国际科技合作基地，1个教育部重点实验室，2个教育部工程研究中心，1个文化部重点实验室，3个省级2011协同创新中心，10个省级重点实验室和工程技术研究中心，2个省级哲学社会科学重点研究基地，5个省级工程实验室，2个省级（技术）研究中心，2个省级国际科技合作基地，1个省级院士专家工作站，1个浙江省新型高校智库，1个浙江省文化厅重点研究基地，1个浙江现代革命历史文化研究基地；拥有3个国家级实验教学示范中心，2个国家级虚拟仿真实验教学示范中心，2个国家工程实践教育中心，1个国家级大学生校外实践教育基地；拥有4门国家精品视频公开课，5门国家精品资源共享课，1门国家级精品在线开放课程，2门国家双语教学示范课，2门教育部来华留学英语授课品牌课程，7本“十二五”国家级规划教材；获国家级教学成果二等奖3项。

学校师资力量雄厚，现有教职工2000余人，其中具有副高以上职称880余人，正高职称270余人。拥有教育部创新团队2个，省重点科技创新团队7个，省重点文化创新团队1个，省高等学校创新团队5个、省高校高水平创新团队3个。国家“万人计划”领军人才4人，教育部长江学者特聘教授1人、讲座教授1人，国家自然科学基金杰出青年科学基金获得者1人，“新世纪百千万人才工程”国家级人选5人，国家有突出贡献中青年专家1人，文化名家暨“四个一批”人才1人，教育部“新世纪优秀人才支持计划”6人；省特级专家3人，省高校“钱江学者”特聘教授12人，省有突出贡献中青年专家9人，省“新世纪151人才工程”第一层次培养人员（含重点资助）31人；享受国务院政府特殊津贴人员31人；聘请了一批国内外知名专家、学者为兼职教授和金沙学者。

学校着力提升科技创新力，在众多领域完成了一系列国家科技计划项目和国家、省部基金科研项目，获得国家技术发明二等奖8项，国家科技进步二等奖7项，国家科技进步三等奖1项，鲁迅文学奖1项，何梁何利基金科学与技术创新奖2项，近五年获得省部级奖励170余项。科技工作综合指标一直稳居浙江省属高校前列。

学校坚持开放办学，已与近30个国家和地区的100余所教育、科研机构建立了合作关系，与毛里求斯大学、苏里南大学共建2所孔子学院，留学生教育迅速发展。

有着120多年办学历史的浙江理工大学，将继续秉承“求知求实、创新创业”的优良传统，大力弘扬“求是笃实”的校风，在新时代把握新机遇，开启新征程，为加快建成位列全国百强、特色鲜明的高水平教学研究型大学而阔步前进！

\end{abstract}

\begin{englishabstract}

Zhejiang Sci-Tech University (ZSTU) is located in Hangzhou, a historical and cultural city, boasts a long history of schooling and covers a campus area of over 2100mu.

Silkworm School, precursor of ZSTU, was founded in 1897 by Lin Qi, governor of Hangzhou prefecture then, with the noble ambition to save the nation by engaging in industry and education, and was one of the earliest new-learning educational institutions ever founded in China. In 1908, Silkworm School was imperialistically promoted by the royal court of the Qing Dynasty to be Advanced Academy of Sericulture and Mulberry owing to its outstanding schooling achievements. The period of turmoil and chaos, from the Revolution of 1911 until the founding of the People’s Republic of China, witnessed the changes of the names and locations of the institution for several times. However, no matter what hardships it underwent, it never gave up its commitment to schooling. After the founding of the PRC in 1949, the academic institution kept forging ahead and composed grand new chapters of its history. In 1959, it began its enrollment of undergraduate candidates. In 1964, the State Council named it Zhejiang Silk Engineering College. In 1979, it started to offer postgraduate programs and was authorized to confer master’s degrees in 1983. With the approval of MOE, it was named Zhejiang Engineering College in 1999 and Zhejiang Sci-Tech University in 2004. In 2006, ZSTU began to have the authority to confer doctoral degrees. In 2010, with the approval of MOE, ZSTU began to be qualified to recommend excellent fresh undergraduates to take master’s degree programs free of examination. In 2017, ZSTU was listed among the Zhejiang Provincial Key Construction Higher Educational Institutions.

ZSTU now is a provincial university focusing on engineering, with distinct characteristics and advantages as well as the coordinated development of the various disciplines of sciences, engineering, arts, economics, management, law, fine art, education, etc. ZSTU consists of 17 affiliated schools or colleges (teaching and research departments) and 1 independent college, with a population of over 26,000 full-time on-campus students, and among whom 3,500 are graduate students. At present, ZSTU offers 63 undergraduate programs of which 8 are the national (MOE) characteristic programs, 1 MOE comprehensive reform trail program, 6 MOE “Outstanding Engineer Education and Cultivation Project” trial programs, 3 “MOE-Autodesk Professional Comprehensive Reform Programs”, 7 “13th five-year” provincially advantageous programs and provincially characteristic programs. ZSTU boasts 2 doctoral degree first-level disciplines (including 12 doctoral degree second-level disciplines), 17 master’s degree first-level disciplines and 6 independently-set master’s degree second-level disciplines (66 master’s degree second-level disciplines, altogether), is authorized to confer ME (including 13 ME categories), MFA, MBA, JM, MAP and MTI, and is authorized to confer master’s degree to candidates of the same educational level. ZSTU has 1 provincial higher education top-key first-level discipline, 3 provincial higher education top-key (first-level) disciplines, 1 provincial higher education key research base of humanities and social sciences, 11 provincial higher education key disciplines, 6 provincial fist-rate disciplines of classification A and 7 provincial first-rate disciplines of classification B. The disciplines of material sciences, chemistry and engineering are listed among the global top 1\% by ESI. ZSTU has 2 state-locality cooperative engineering laboratories, 2 national international sci-tech cooperation bases, 1 MOE key laboratory, 2 MOE engineering research centers, 1 MOC key laboratory, 3 provincial 2011 coordinated innovation centers, 10 provincial key laboratories and engineering technology research centers, 1 provincial philosophy and social sciences key research base, 5 provincial engineering laboratories, 2 provincial (technological) research centers, 2 provincial international sci-tech cooperation bases, 1 provincial academician and expert work station, 1 provincial DOC key research base, 1 provincial modern revolution history cultural research base, 3 national experimental teaching demonstration centers, 2 national virtual simulation experimental teaching demonstration centers, 2 national engineering practical education centers and 1 national college student off-campus practical education base. ZSTU has 4 national choicest video public courses, 5 national choicest resource sharing courses, 2 national bilingual teaching demonstration courses, 2 MOE foreign student in-English brand courses and 7 “12th five-year” nationally-planned textbooks. ZSTU has achieved 3 second prizes of national teaching awards.

ZSTU boasts a high-level faculty and staff of more than 2,000 personages, with over 880 associate professors and 260 professors. It has 2 MOE innovation teams, 7 provincial key sci-tech innovation teams, 1 provincial key cultural innovation team, 5 provincial higher-education innovation teams and 2 post-doctor scientific research mobile stations. ZSTU has 3 national “Ten Thousand Talent Program” leading talents, 1 MOE Yangtze River distinguished professor, 1 MOE Yangtze River lecturing professor, 1 winner of national NSFC outstanding young scientist foundation, 3 national “One Thousand Talent Program” experts (including 1 national “One Thousand Foreign Expert Program” expert), 5 national “New-century Thousands-of-talent Program” talents, 1 national outstanding contribution middle-young-aged expert, 1 Cultural Master and “Four Batches” talent, 5 MOE “New-century Excellent Talent Support Program” talents, 1 provincial top expert, 5 provincial “One Thousand Talent Program” (including 2 Seagull Program talents), 10 provincial higher education “Qiantangjiang River Scholar” distinguished professors, 9 provincial outstanding contribution middle-young-aged experts, 27 first-level cultivation talents of provincial “New-century 151 Talent Program”, and 31 State Council government special subsidy personages. Additionally, ZSTU employs massive renowned experts and scholars from home and abroad as part-time professors and Jinsha Scholars.

ZSTU has ever been attaching importance to scientific research work, the research strength has been progressively enhanced, and ZSTU has successfully undertaken series of national sci-tech planning projects, as well as other national or provincial research projects in various fields. ZSTU has won 7 second prizes of State Technological Invention Award, 5 second prizes of State Sci-tech Progress Award, 1 Lu Xun Literature Award, 1 Ho Leung Ho Lee Foundation sci-tech innovation prize and 150 or so provincial or ministry prizes. The composite index of the sci-tech work has been ranking among the top of the provincial universities and colleges.

ZSTU has been sticking to its open schooling, has established its cooperation with over 100 educational and sci-tech research institutions of over 20 countries and areas, has cooperatively established 2 Confucius Institutes with University of Mauritius and University of Suriname, and the foreign student education has been rapidly developing.

ZSTU with a schooling history of 120 years is sparing no efforts to carry forward its school spirit of “pursuing truth and uprightness”, to take new opportunities and start new journeys in the new era, and to stride forward to build the university into a high-level teaching-research-style institution of distinct characteristics as soon as possible.

\end{englishabstract}



% 目录 部分课程论文需要，可以自行选择添加或者去除
% \tableofcontents

\else
  \ifzstu@submit\relax
    \includepdf{pdf/original.pdf}
    \cleardoublepage
%    \includepdf{pdf/authorization.pdf}
%    \cleardoublepage
  \else
    \ifzstu@review\relax
    % exclude the original claim and authorization
    \else
      \makeDeclareOriginal
%      \makeDeclareAuthorization
    \fi
  \fi
  \frontmatter % 使用罗马数字对前言编号

  % 摘要
  %# -*- coding: utf-8-unix -*-
% !TEX program = xelatex
% !TEX root = ../thesis.tex
% !TEX encoding = UTF-8 Unicode
%%==================================================
%% abstract.tex for SJTU Master Thesis
%%==================================================

\begin{abstract}

学校办学历史悠久。浙江理工大学坐落在历史文化名城杭州市，占地面积2100余亩。学校前身——蚕学馆，是杭州知府林启为实现实业救国、教育救国的宏愿于1897年创办的，是我国最早创办的新学教育机构之一。1908年，因办学成绩卓著，被清政府御批升格为“高等蚕桑学堂”。辛亥革命至解放前夕，因时局动乱，学校几度易名、数迁校址，在风雨沧桑中生生不息、坚持办学。新中国成立后，学校不断开拓进取，绘就了新的历史篇章。学校1959年开始招收本科生，1964年由国务院定名为浙江丝绸工学院，1979年开始招收硕士研究生，1983年获硕士学位授予权。1999年，经教育部批准，学校更名为浙江工程学院。2004年，经教育部批准，学校更名为浙江理工大学。2006年，获得博士学位授予权。2010年，经教育部批准，获得推荐优秀应届本科毕业生免试攻读硕士学位研究生资格。2017年，学校入选浙江省重点建设高校。

今天的浙江理工大学是一所以工为主，特色鲜明，优势突出，理、工、文、经、管、法、艺术、教育等多学科协调发展的省重点建设高校。学校下设17个学院（教研部），举办1所独立学院，现有全日制在校学生27000余人，其中研究生3800余人。学校现有本科专业65个，其中国家（教育部）特色专业8个，教育部综合改革试点专业1个、教育部“卓越工程师教育培养计划”试点专业6个、“教育部-欧特克公司专业综合改革项目”3个，“十三五”省优势专业7个、省特色专业5个。拥有省重点建设高校优势特色学科2个，博士学位授权一级学科3个、硕士学位授权一级学科26个，并可授予专业学位类别10种，博士后科研流动站2个。拥有省高校重中之重一级学科1个、省高校重中之重（一级）学科3个、省高校人文社会科学重点研究基地1个、省高校重点学科11个，省一流学科A类6个、一流学科B类7个，材料科学、化学、工程学3个学科进入ESI全球排名前1\%；拥有2个国家地方联合工程实验室，1个国家地方联合工程研究中心，3个国家国际科技合作基地，1个教育部重点实验室，2个教育部工程研究中心，1个文化部重点实验室，3个省级2011协同创新中心，10个省级重点实验室和工程技术研究中心，2个省级哲学社会科学重点研究基地，5个省级工程实验室，2个省级（技术）研究中心，2个省级国际科技合作基地，1个省级院士专家工作站，1个浙江省新型高校智库，1个浙江省文化厅重点研究基地，1个浙江现代革命历史文化研究基地；拥有3个国家级实验教学示范中心，2个国家级虚拟仿真实验教学示范中心，2个国家工程实践教育中心，1个国家级大学生校外实践教育基地；拥有4门国家精品视频公开课，5门国家精品资源共享课，1门国家级精品在线开放课程，2门国家双语教学示范课，2门教育部来华留学英语授课品牌课程，7本“十二五”国家级规划教材；获国家级教学成果二等奖3项。

学校师资力量雄厚，现有教职工2000余人，其中具有副高以上职称880余人，正高职称270余人。拥有教育部创新团队2个，省重点科技创新团队7个，省重点文化创新团队1个，省高等学校创新团队5个、省高校高水平创新团队3个。国家“万人计划”领军人才4人，教育部长江学者特聘教授1人、讲座教授1人，国家自然科学基金杰出青年科学基金获得者1人，“新世纪百千万人才工程”国家级人选5人，国家有突出贡献中青年专家1人，文化名家暨“四个一批”人才1人，教育部“新世纪优秀人才支持计划”6人；省特级专家3人，省高校“钱江学者”特聘教授12人，省有突出贡献中青年专家9人，省“新世纪151人才工程”第一层次培养人员（含重点资助）31人；享受国务院政府特殊津贴人员31人；聘请了一批国内外知名专家、学者为兼职教授和金沙学者。

学校着力提升科技创新力，在众多领域完成了一系列国家科技计划项目和国家、省部基金科研项目，获得国家技术发明二等奖8项，国家科技进步二等奖7项，国家科技进步三等奖1项，鲁迅文学奖1项，何梁何利基金科学与技术创新奖2项，近五年获得省部级奖励170余项。科技工作综合指标一直稳居浙江省属高校前列。

学校坚持开放办学，已与近30个国家和地区的100余所教育、科研机构建立了合作关系，与毛里求斯大学、苏里南大学共建2所孔子学院，留学生教育迅速发展。

有着120多年办学历史的浙江理工大学，将继续秉承“求知求实、创新创业”的优良传统，大力弘扬“求是笃实”的校风，在新时代把握新机遇，开启新征程，为加快建成位列全国百强、特色鲜明的高水平教学研究型大学而阔步前进！

\end{abstract}

\begin{englishabstract}

Zhejiang Sci-Tech University (ZSTU) is located in Hangzhou, a historical and cultural city, boasts a long history of schooling and covers a campus area of over 2100mu.

Silkworm School, precursor of ZSTU, was founded in 1897 by Lin Qi, governor of Hangzhou prefecture then, with the noble ambition to save the nation by engaging in industry and education, and was one of the earliest new-learning educational institutions ever founded in China. In 1908, Silkworm School was imperialistically promoted by the royal court of the Qing Dynasty to be Advanced Academy of Sericulture and Mulberry owing to its outstanding schooling achievements. The period of turmoil and chaos, from the Revolution of 1911 until the founding of the People’s Republic of China, witnessed the changes of the names and locations of the institution for several times. However, no matter what hardships it underwent, it never gave up its commitment to schooling. After the founding of the PRC in 1949, the academic institution kept forging ahead and composed grand new chapters of its history. In 1959, it began its enrollment of undergraduate candidates. In 1964, the State Council named it Zhejiang Silk Engineering College. In 1979, it started to offer postgraduate programs and was authorized to confer master’s degrees in 1983. With the approval of MOE, it was named Zhejiang Engineering College in 1999 and Zhejiang Sci-Tech University in 2004. In 2006, ZSTU began to have the authority to confer doctoral degrees. In 2010, with the approval of MOE, ZSTU began to be qualified to recommend excellent fresh undergraduates to take master’s degree programs free of examination. In 2017, ZSTU was listed among the Zhejiang Provincial Key Construction Higher Educational Institutions.

ZSTU now is a provincial university focusing on engineering, with distinct characteristics and advantages as well as the coordinated development of the various disciplines of sciences, engineering, arts, economics, management, law, fine art, education, etc. ZSTU consists of 17 affiliated schools or colleges (teaching and research departments) and 1 independent college, with a population of over 26,000 full-time on-campus students, and among whom 3,500 are graduate students. At present, ZSTU offers 63 undergraduate programs of which 8 are the national (MOE) characteristic programs, 1 MOE comprehensive reform trail program, 6 MOE “Outstanding Engineer Education and Cultivation Project” trial programs, 3 “MOE-Autodesk Professional Comprehensive Reform Programs”, 7 “13th five-year” provincially advantageous programs and provincially characteristic programs. ZSTU boasts 2 doctoral degree first-level disciplines (including 12 doctoral degree second-level disciplines), 17 master’s degree first-level disciplines and 6 independently-set master’s degree second-level disciplines (66 master’s degree second-level disciplines, altogether), is authorized to confer ME (including 13 ME categories), MFA, MBA, JM, MAP and MTI, and is authorized to confer master’s degree to candidates of the same educational level. ZSTU has 1 provincial higher education top-key first-level discipline, 3 provincial higher education top-key (first-level) disciplines, 1 provincial higher education key research base of humanities and social sciences, 11 provincial higher education key disciplines, 6 provincial fist-rate disciplines of classification A and 7 provincial first-rate disciplines of classification B. The disciplines of material sciences, chemistry and engineering are listed among the global top 1\% by ESI. ZSTU has 2 state-locality cooperative engineering laboratories, 2 national international sci-tech cooperation bases, 1 MOE key laboratory, 2 MOE engineering research centers, 1 MOC key laboratory, 3 provincial 2011 coordinated innovation centers, 10 provincial key laboratories and engineering technology research centers, 1 provincial philosophy and social sciences key research base, 5 provincial engineering laboratories, 2 provincial (technological) research centers, 2 provincial international sci-tech cooperation bases, 1 provincial academician and expert work station, 1 provincial DOC key research base, 1 provincial modern revolution history cultural research base, 3 national experimental teaching demonstration centers, 2 national virtual simulation experimental teaching demonstration centers, 2 national engineering practical education centers and 1 national college student off-campus practical education base. ZSTU has 4 national choicest video public courses, 5 national choicest resource sharing courses, 2 national bilingual teaching demonstration courses, 2 MOE foreign student in-English brand courses and 7 “12th five-year” nationally-planned textbooks. ZSTU has achieved 3 second prizes of national teaching awards.

ZSTU boasts a high-level faculty and staff of more than 2,000 personages, with over 880 associate professors and 260 professors. It has 2 MOE innovation teams, 7 provincial key sci-tech innovation teams, 1 provincial key cultural innovation team, 5 provincial higher-education innovation teams and 2 post-doctor scientific research mobile stations. ZSTU has 3 national “Ten Thousand Talent Program” leading talents, 1 MOE Yangtze River distinguished professor, 1 MOE Yangtze River lecturing professor, 1 winner of national NSFC outstanding young scientist foundation, 3 national “One Thousand Talent Program” experts (including 1 national “One Thousand Foreign Expert Program” expert), 5 national “New-century Thousands-of-talent Program” talents, 1 national outstanding contribution middle-young-aged expert, 1 Cultural Master and “Four Batches” talent, 5 MOE “New-century Excellent Talent Support Program” talents, 1 provincial top expert, 5 provincial “One Thousand Talent Program” (including 2 Seagull Program talents), 10 provincial higher education “Qiantangjiang River Scholar” distinguished professors, 9 provincial outstanding contribution middle-young-aged experts, 27 first-level cultivation talents of provincial “New-century 151 Talent Program”, and 31 State Council government special subsidy personages. Additionally, ZSTU employs massive renowned experts and scholars from home and abroad as part-time professors and Jinsha Scholars.

ZSTU has ever been attaching importance to scientific research work, the research strength has been progressively enhanced, and ZSTU has successfully undertaken series of national sci-tech planning projects, as well as other national or provincial research projects in various fields. ZSTU has won 7 second prizes of State Technological Invention Award, 5 second prizes of State Sci-tech Progress Award, 1 Lu Xun Literature Award, 1 Ho Leung Ho Lee Foundation sci-tech innovation prize and 150 or so provincial or ministry prizes. The composite index of the sci-tech work has been ranking among the top of the provincial universities and colleges.

ZSTU has been sticking to its open schooling, has established its cooperation with over 100 educational and sci-tech research institutions of over 20 countries and areas, has cooperatively established 2 Confucius Institutes with University of Mauritius and University of Suriname, and the foreign student education has been rapidly developing.

ZSTU with a schooling history of 120 years is sparing no efforts to carry forward its school spirit of “pursuing truth and uprightness”, to take new opportunities and start new journeys in the new era, and to stride forward to build the university into a high-level teaching-research-style institution of distinct characteristics as soon as possible.

\end{englishabstract}



  % 目录、插图目录、表格目录
  \tableofcontents
  \listoffigures
  \addcontentsline{toc}{chapter}{\listfigurename}     % 将插图目录加入全文目录
  \listoftables
  \addcontentsline{toc}{chapter}{\listtablename}      % 将表格目录加入全文目录
  \listofalgorithms
  \addcontentsline{toc}{chapter}{\listalgorithmname}  % 将算法目录加入全文目录

  \include{tex/symbol} % 主要符号、缩略词对照表
\fi

\makeatother
\mainmatter % 使用阿拉伯数字对正文编号

% 正文内容
%# -*- coding: utf-8-unix -*-
% !TEX program = xelatex
% !TEX root = ../thesis.tex
% !TEX encoding = UTF-8 Unicode
%%==================================================
%% chapter01.tex for ZSTU Master Thesis
%%==================================================

%\bibliographystyle{zstu2}%[此处用于每章都生产参考文献]
\chapter{这是什么}
\label{chap:intro}

这是上海交通大学(非官方)学位论文 \LaTeX 模板，当前版本是 \version 。

最早的一版学位模板是一位热心的物理系同学制作的。
那份模板参考了自动化所学位论文模板，使用了CASthesis.cls文档类，中文字符处理则采用当时最为流行的 \CJKLaTeX 方案。
我根据交大研究生院对学位论文的要求
\footnote{\url{http://www.gs.zstu.edu.cn/policy/fileShow.ahtml?id=130}}
，结合少量个人审美喜好，完成了一份基本可用的交大 \LaTeX 学位论文模板。
但是，搭建一个 \CJKLaTeX 环境并不简单，单单在Linux下配置环境和添加中文字体，就足够让新手打退堂鼓。
在William Wang的建议下，我开始着手把模板向 \XeTeX 引擎移植。
他完成了最初的移植，多亏了他出色的工作，后续的改善工作也得以顺利进行。

随着我对 \LaTeX 系统认知增加，我又断断续续做了一些完善模板的工作，在原有硕士学位论文模板的基础上完成了交大学士和博士学位论文模板。

现在，交大学位论文模板ZSTUTHesis代码在github
\footnote{\url{https://github.com/sjtug/SJTUThesis}}
上维护。
你可以\href{https://github.com/sjtug/SJTUThesis/issues}{在github上开issue}
、或者在\href{https://bbs.sjtu.edu.cn/bbsdoc?board=TeX_LaTeX}{水源LaTeX版}发帖来反映遇到的问题。

\section{使用模板}

\subsection{准备工作}
\label{sec:requirements}

要使用这个模板撰写学位论文，需要在\emph{TeX系统}、\emph{TeX技能}上有所准备。

\begin{itemize}[noitemsep,topsep=0pt,parsep=0pt,partopsep=0pt]
	\item {\TeX}系统：所使用的{\TeX}系统要支持 \XeTeX 引擎，且带有ctex 2.x宏包，以2017年或更新版本的\emph{完整}TeXLive、MacTeX发行版为佳。
	\item TeX技能：尽管提供了对模板的必要说明，但这不是一份“ \LaTeX 入门文档”。在使用前请先通读其他入门文档。
	\item 针对Windows用户的额外需求：学位论文模本分别使用git和GNUMake进行版本控制和构建，建议从Cygwin\footnote{\url{http://cygwin.com}}安装这两个工具。
\end{itemize}

\subsection{模板选项}
\label{sec:thesisoption}

sjtuthesis提供了一些常用选项，在thesis.tex在导入sjtuthesis模板类时，可以组合使用。
这些选项包括：

\begin{itemize}[noitemsep,topsep=0pt,parsep=0pt,partopsep=0pt]
	\item 学位类型：bachelor(学位)、master(硕士)、doctor(博士)，是必选项。
	\item 中文字体：fandol(Fandol 开源字体)、windows(Windows 系统下的中文字体)、mac(macOS 系统下的华文字体)、ubuntu(Ubuntu 系统下的文泉驿和文鼎字体)、adobe(Adobe 公司的中文字体)、founder(方正公司的中文字体)，默认根据操作系统自动配置。
	\item 英文模版：使用english选项启用英文模版。
	\item 盲审选项：使用review选项后，论文作者、学号、导师姓名、致谢、发表论文和参与项目将被隐去。
\end{itemize}

\subsection{编译模板}
\label{sec:process}

模板默认使用GNUMake构建，GNUMake将调用latemk工具自动完成模板多轮编译：

\begin{lstlisting}[basicstyle=\small\ttfamily, caption={编译模板}, numbers=none]
make clean thesis.pdf
\end{lstlisting}

若需要生成包含“原创性声明扫描件”的学位论文文档，请将扫描件保存为statement.pdf，然后调用make生成submit.pdf。

\begin{lstlisting}[basicstyle=\small\ttfamily, caption={生成用于提交的学位论文}, numbers=none]
make clean submit.pdf
\end{lstlisting}

编译失败时，可以尝试手动逐次编译，定位故障。

\begin{lstlisting}[basicstyle=\small\ttfamily, caption={手动逐次编译}, numbers=none]
xelatex -no-pdf thesis
biber --debug thesis
xelatex thesis
xelatex thesis
\end{lstlisting}

\subsection{模板文件布局}
\label{sec:layout}

\begin{lstlisting}[basicstyle=\small\ttfamily,caption={模板文件布局},label=layout,float,numbers=none]
├── LICENSE
├── Makefile
├── README.md
├── bib
│   ├── chap1.bib
│   └── chap2.bib
├── bst
│   └── GBT7714-2005NLang.bst
├── figure
│   ├── chap2
│   │   ├── sjtulogo.eps
│   │   ├── sjtulogo.jpg
│   │   ├── sjtulogo.pdf
│   │   └── sjtulogo.png
│   └── sjtubanner.png
├── sjtuthesis.cfg
├── sjtuthesis.cls
├── statement.pdf
├── submit.pdf
├── tex
│   ├── abstract.tex
│   ├── ack.tex
│   ├── app_cjk.tex
│   ├── app_eq.tex
│   ├── app_log.tex
│   ├── chapter01.tex
│   ├── chapter02.tex
│   ├── chapter03.tex
│   ├── conclusion.tex
│   ├── id.tex
│   ├── patents.tex
│   ├── projects.tex
│   ├── pub.tex
│   └── symbol.tex
└── thesis.tex
\end{lstlisting}

本节介绍学位论文模板中木要文件和目录的功能。

\subsubsection{格式控制文件}
\label{sec:format}

格式控制文件控制着论文的表现形式，包括sjtuthesis.cfg和sjtuthesis.cls。
其中，“cls”控制论文主体格式，“cfg”为配置文件。

\subsubsection{主控文件thesis.tex}
\label{sec:thesistex}

主控文件thesis.tex的作用就是将你分散在多个文件中的内容“整合”成一篇完整的论文。
使用这个模板撰写学位论文时，你的学位论文内容和素材会被“拆散”到各个文件中：
譬如各章正文、各个附录、各章参考文献等等。
在thesis.tex中通过“include”命令将论文的各个部分包含进来，从而形成一篇结构完成的论文。
对模板定制时引入的宏包，建议放在导言区。

\subsubsection{各章源文件tex}
\label{sec:thesisbody}

这一部分是论文的主体，是以“章”为单位划分的，包括：

\begin{itemize}[noitemsep,topsep=0pt,parsep=0pt,partopsep=0pt]
	\item 中英文摘要(abstract.tex)。前言(frontmatter)的其他部分，中英文封面、原创性声明、授权信息在sjtuthesis.cls中定义，不单独分离为tex文件。
不单独弄成文件。
	\item 正文(mainmatter)——学位论文正文的各章内容，源文件是chapter\emph{xxx}.tex。
	\item 附录(app\emph{xx}.tex)、致谢(ack.tex)、攻读学位论文期间发表的学术论文目录(pub.tex)、个人简历(resume.tex)组成正文后的部分(backmatter)。
参考文献列表由bibtex插入，不作为一个单独的文件。
\end{itemize}

\subsubsection{图片文件夹figure}
\label{sec:fig}

figure文件夹放置了需要插入文档中的图片文件(支持PNG/JPG/PDF/EPS格式的图片)，可以在按照章节划分子目录。
模板文件中使用\verb|\graphicspath|命令定义了图片存储的顶层目录，在插入图片时，顶层目录名“figure”可省略。

\subsubsection{参考文献数据库bib}
\label{sec:bib}

目前参考文件数据库目录只存放一个参考文件数据库thesis.bib。
关于参考文献引用，可参考第\ref{chap:example}章中的例子。


%# -*- coding: utf-8-unix -*-
% !TEX program = xelatex
% !TEX root = ../thesis.tex
% !TEX encoding = UTF-8 Unicode
%%==================================================
%% chapter02.tex for ZSTU Master Thesis
%% based on CASthesis
%% modified by wei.jianwen@gmail.com
%% Encoding: UTF-8
%%==================================================

\chapter{{\LaTeX} 排版例子}
\label{chap:example}

\section{列表环境}
\label{sec:list}

\subsection{无序列表}
\label{sec:unorderlist}

以下是一个无序列表的例子，列表的每个条目单独分段。

\begin{itemize}
  \item 这是一个无序列表。
  \item 这是一个无序列表。
  \item 这是一个无序列表。
\end{itemize}

使用\verb+itemize*+环境可以创建行内无序列表。
\begin{itemize*}
  \item 这是一个无序列表。
  \item 这是一个无序列表。
  \item 这是一个无序列表。
\end{itemize*}
行内无序列表条目不单独分段，所有内容直接插入在原文的段落中。

\subsection{有序列表}
\label{sec:orderlist}

使用环境\verb+enumerate+和\verb+enumerate*+创建有序列表,
使用方法无序列表类似。

\begin{enumerate}
  \item 这是一个有序列表。
  \item 这是一个有序列表。
  \item 这是一个有序列表。
\end{enumerate}

使用\verb+enumerate*+环境可以创建行内有序列表。
\begin{enumerate*}
  \item 这是一个默认有序列表。
  \item 这是一个默认有序列表。
  \item 这是一个默认有序列表。
\end{enumerate*}
行内有序列表条目不单独分段，所有内容直接插入在原文的段落中。

\subsection{描述型列表}

使用环境\verb+description+可创建带有主题词的列表，条目语法是\verb+\item[主题] 内容+。
\begin{description}
    \item[主题一] 详细内容
    \item[主题二] 详细内容
    \item[主题三] 详细内容 \ldots
\end{description}

\subsection{自定义列表样式}

可以使用\verb+label+参数控制列表的样式，
详细可以参考WikiBooks\footnote{\url{https://en.wikibooks.org/wiki/LaTeX/List_Structures\#Customizing_lists}}。
比如一个自定义样式的行内有序列表
\begin{enumerate*}[label=\itshape\alph*)\upshape]
  \item 这是一个自定义样式有序列表。
  \item 这是一个自定义样式有序列表。
  \item 这是一个自定义样式有序列表。
\end{enumerate*}

\section{数学排版}
\label{sec:matheq}

\subsection{公式排版}
\label{sec:eqformat}

这里有举一个长公式排版的例子，来自\href{http://www.tex.ac.uk/tex-archive/info/math/voss/mathmode/Mathmode.pdf}{《Math mode》}：

\begin {multline}
  \frac {1}{2}\Delta (f_{ij}f^{ij})=
  2\left (\sum _{i<j}\chi _{ij}(\sigma _{i}-
    \sigma _{j}) ^{2}+ f^{ij}\nabla _{j}\nabla _{i}(\Delta f)+\right .\\
  \left .+\nabla _{k}f_{ij}\nabla ^{k}f^{ij}+
    f^{ij}f^{k}\left [2\nabla _{i}R_{jk}-
      \nabla _{k}R_{ij}\right ]\vphantom {\sum _{i<j}}\right )
\end{multline}

\subsection{SI单位}

使用\verb+siunitx+宏包可以方便地输入SI单位制单位，例如\verb+\SI{5}{\um}+可以得到\SI{5}{\um}。

\subsubsection{一个四级标题}
\label{sec:depth4}

这是全文唯一的一个四级标题。在这部分中将演示了mathtools宏包中可伸长符号（箭头、等号的例子）的例子。

\begin{displaymath}
    A \xleftarrow[n=0]{} B \xrightarrow[LongLongLongLong]{n>0} C
\end{displaymath}

\begin{eqnarray}
  f(x) & \xleftrightarrow[]{A=B}  & B \\
  & \xleftharpoondown[below]{above} & B \nonumber \\
  & \xLeftrightarrow[below]{above} & B
\end{eqnarray}

又如：

\begin{align}
  \label{eq:none}
  & I(X_3;X_4)-I(X_3;X_4\mid{}X_1)-I(X_3;X_4\mid{}X_2) \nonumber \\
  = & [I(X_3;X_4)-I(X_3;X_4\mid{}X_1)]-I(X_3;X_4\mid{}\tilde{X}_2) \\
  = & I(X_1;X_3;X_4)-I(X_3;X_4\mid{}\tilde{X}_2)
\end{align}

\subsection{定理环境}

模板中定义了丰富的定理环境
algo(算法)，thm(定理)，lem(引理)，prop(命题)，cor(推论)，defn(定义)，conj(猜想)，exmp(例)，rem(注)，case(情形)，
bthm(断言定理)，blem(断言引理)，bprop(断言命题)，bcor(断言推论)。
amsmath还提供了一个proof(证明)的环境。
这里举一个“定理”和“证明”的例子。
\begin{thm}[留数定理]
\label{thm:res}
  假设$U$是复平面上的一个单连通开子集，$a_1,\ldots,a_n$是复平面上有限个点，$f$是定义在$U\backslash \{a_1,\ldots,a_n\}$上的全纯函数，
  如果$\gamma$是一条把$a_1,\ldots,a_n$包围起来的可求长曲线，但不经过任何一个$a_k$，并且其起点与终点重合，那么：

  \begin{equation}
    \label{eq:res}
    \ointop_{\gamma}f(z)\,\mathrm{d}z = 2\uppi\mathbf{i}\sum^n_{k=1}\mathrm{I}(\gamma,a_k)\mathrm{Res}(f,a_k)
  \end{equation}

  如果$\gamma$是若尔当曲线，那么$\mathrm{I}(\gamma, a_k)=1$，因此：

  \begin{equation}
    \label{eq:resthm}
    \ointop_{\gamma}f(z)\,\mathrm{d}z = 2\uppi\mathbf{i}\sum^n_{k=1}\mathrm{Res}(f,a_k)
  \end{equation}

      % \oint_\gamma f(z)\, dz = 2\pi i \sum_{k=1}^n \mathrm{Res}(f, a_k ).

  在这里，$\mathrm{Res}(f, a_k)$表示$f$在点$a_k$的留数，$\mathrm{I}(\gamma,a_k)$表示$\gamma$关于点$a_k$的卷绕数。
  卷绕数是一个整数，它描述了曲线$\gamma$绕过点$a_k$的次数。如果$\gamma$依逆时针方向绕着$a_k$移动，卷绕数就是一个正数，
  如果$\gamma$根本不绕过$a_k$，卷绕数就是零。

  定理\ref{thm:res}的证明。

  \begin{proof}
    首先，由……

    其次，……

    所以……
  \end{proof}
\end{thm}

上面的公式例子中，有一些细节希望大家注意。微分号d应该使用“直立体”也就是用mathrm包围起来。
并且，微分号和被积函数之间应该有一段小间隔，可以插入\verb+\,+得到。
斜体的$d$通常只作为一般变量。
i,j作为虚数单位时，也应该使用“直立体”为了明显，还加上了粗体，例如\verb+\mathbf{i}+。斜体$i,j$通常用作表示“序号”。
其他字母在表示常量时，也推荐使用“直立体”譬如，圆周率$\uppi$（需要upgreek宏包），自然对数的底$\mathrm{e}$。
不过，我个人觉得斜体的$e$和$\pi$很潇洒，在不至于引起混淆的情况下，我也用这两个字母的斜体表示对应的常量。


\section{向文档中插入图像}
\label{sec:insertimage}

\subsection{支持的图片格式}
\label{sec:imageformat}

\XeTeX 可以很方便地插入PDF、PNG、JPG格式的图片。

插入PNG/JPG的例子如\ref{fig:SRR}所示。
这两个水平并列放置的图共享一个“图标题”(table caption)，没有各自的小标题。

\begin{figure}[!htp]
  \centering
  \includegraphics[width=4cm]{example/zstulogo.png}
  \hspace{1cm}
  \includegraphics[width=4cm]{example/zstulogo.jpg}
  \bicaption[这里将出现在插图索引中]
    {中文题图}
    {English caption}
  \label{fig:SRR}
\end{figure}

这里还有插入EPS图像和PDF图像的例子，如图\ref{fig:epspdf:a}和图\ref{fig:epspdf:b}。这里将EPS和PDF图片作为子图插入，每个子图有自己的小标题。子图标题使用subcaption宏包添加。

\begin{figure}[!htp]
  \centering
  \subcaptionbox{EPS 图像\label{fig:epspdf:a}}[3cm] %标题的长度，超过则会换行，如下一个小图。
    {\includegraphics[height=2.5cm]{example/zstulogo.eps}}
  \hspace{4em}
  \subcaptionbox{PDF 图像，注意这个图略矮些。如果标题很长的话，它会自动换行\label{fig:epspdf:b}}
    {\includegraphics[height=2cm]{zstulogo.pdf}}
  \bicaption{插入eps和pdf的例子（使用 subcaptionbox 方式）}{An EPS and PDF demo with subcaptionbox}
  \label{fig:pdfeps-subcaptionbox}
\end{figure}

\begin{figure}[!htp]
  \centering
  \begin{subfigure}{2.5cm}
    \centering
    \includegraphics[height=2.5cm]{example/zstulogo.eps}
    \caption{EPS 图像}
  \end{subfigure}
  \hspace{4em}
  \begin{subfigure}{0.4\textwidth}
    \centering
    \includegraphics[height=2cm]{zstulogo.pdf}
    \caption{PDF 图像，注意这个图略矮些。subfigure中同一行的子图在顶端对齐。}
  \end{subfigure}
  \bicaption{插入eps和pdf的例子（使用 subfigure 方式）}{An EPS and PDF demo with subfigure}
  \label{fig:pdfeps-subfigure}
\end{figure}

更多关于 \LaTeX 插图的例子可以参考\href{http://www.cs.duke.edu/junhu/Graphics3.pdf}{《\LaTeX 插图指南》}。

\subsection{长标题的换行}
\label{sec:longcaption}

图\ref{fig:longcaptionbad}和图\ref{fig:longcaptiongood}都有比较长图标题，通过对比发现，图\ref{fig:longcaptiongood}的换行效果更好一些。
其中使用了minipage环境来限制整个浮动体的宽度。

\begin{figure}[!htp]
  \centering
  \includegraphics[width=4cm]{zstubadge.pdf}
  \bicaption[这里将出现在插图索引]
    {上海交通大学是我国历史最悠久的高等学府之一，是教育部直属、教育部与上海市共建的全国重点大学.}
    {Where there is a will, there is a way.}
 \label{fig:longcaptionbad}
\end{figure}

\begin{figure}[!htbp]
  \centering
  \begin{minipage}[b]{0.6\textwidth}
    \centering
    \includegraphics[width=4cm]{zstubadge.pdf}
    \bicaption[出现在插图索引中]
      {上海交通大学是我国历史最悠久的高等学府之一，是教育部直属、教育部与上海市共建的全国重点大学.}
      {Where there is a will, there is a way.}
    \label{fig:longcaptiongood}
  \end{minipage}
\end{figure}

\subsection{添加图注}

当插图中组成部件由数字或字母等编号表示时，可在插图下方添加图注进行说明，如图\ref{fig:cn_100t}所示。

\begin{figure}[!htp]
  \centering
  \includegraphics[width=0.3\textwidth]{example/cn_100t.png}\
  \begin{center}
    \small\kaishu 1.立柱 2.提升释放机构 3.标准冲击加速度计 \\ 4.导轨 5.重锤 6.被校力传感器 7.底座
  \end{center}
  \vspace{-1em}
  \bicaption[出现在插图索引中]
    {示例图片来源于\parencite{he1999}}
    {Stay hungry, stay foolish.}
 \label{fig:cn_100t}
\end{figure}

\subsection{绘制流程图}

图\ref{fig:flow_chart}是一张流程图示意。使用tikz环境，搭配四种预定义节点（\verb+startstop+、\verb+process+、\verb+decision+和\verb+io+），可以容易地绘制出流程图。
\begin{figure}[!htp]
    \centering
    \resizebox{6cm}{!}{\input{figure/example/flow_chart.tex}}
    \bicaption{绘制流程图效果}{Flow chart}
    \label{fig:flow_chart}
\end{figure}

\clearpage

\section{表格}
\label{sec:tab}

这一节给出的是一些表格的例子，如表\ref{tab:firstone}所示。

\begin{table}[!hpb]
  \centering
  \bicaption[指向一个表格的表目录索引]
    {一个颇为标准的三线表格\footnotemark[1]}
    {A Table}
  \label{tab:firstone}
  \begin{tabular}{@{}llr@{}} \toprule
    \multicolumn{2}{c}{Item} \\ \cmidrule(r){1-2}
    Animal & Description & Price (\$)\\ \midrule
    Gnat & per gram & 13.65 \\
    & each & 0.01 \\
    Gnu & stuffed & 92.50 \\
    Emu & stuffed & 33.33 \\
    Armadillo & frozen & 8.99 \\ \bottomrule
  \end{tabular}
\end{table}
\footnotetext[1]{这个例子来自\href{http://www.ctan.org/tex-archive/macros/latex/contrib/booktabs/booktabs.pdf}{《Publication quality tables in LATEX》}(booktabs宏包的文档)。这也是一个在表格中使用脚注的例子，请留意与threeparttable实现的效果有何不同。}

下面一个是一个更复杂的表格，用threeparttable实现带有脚注的表格，如表\ref{tab:footnote}。

\begin{table}[!htpb]
  \bicaption[出现在表目录的标题]
    {一个带有脚注的表格的例子}
    {A Table with footnotes}
  \label{tab:footnote}
  \centering
  \begin{threeparttable}[b]
     \begin{tabular}{ccd{4}cccc}
      \toprule
      \multirow{2}{6mm}{total}&\multicolumn{2}{c}{20\tnote{1}} & \multicolumn{2}{c}{40} &  \multicolumn{2}{c}{60}\\
      \cmidrule(lr){2-3}\cmidrule(lr){4-5}\cmidrule(lr){6-7}
      &www & \multicolumn{1}{c}{k} & www & k & www & k \\ % 使用说明符 d 的列会自动进入数学模式，使用 \multicolumn 对文字表头做特殊处理
      \midrule
      &$\underset{(2.12)}{4.22}$ & 120.0140\tnote{2} & 333.15 & 0.0411 & 444.99 & 0.1387 \\
      &168.6123 & 10.86 & 255.37 & 0.0353 & 376.14 & 0.1058 \\
      &6.761    & 0.007 & 235.37 & 0.0267 & 348.66 & 0.1010 \\
      \bottomrule
    \end{tabular}
    \begin{tablenotes}
    \item [1] the first note.% or \item [a]
    \item [2] the second note.% or \item [b]
    \end{tablenotes}
  \end{threeparttable}
\end{table}

\section{参考文献管理}

 \LaTeX 具有将参考文献内容和表现形式分开管理的能力，涉及三个要素：参考文献数据库、参考文献引用格式、在正文中引用参考文献。
这样的流程需要多次编译：

\begin{enumerate}[noitemsep,topsep=0pt,parsep=0pt,partopsep=0pt]
	\item 用户将论文中需要引用的参考文献条目，录入纯文本数据库文件(bib文件)。
	\item 调用xelatex对论文模板做第一次编译，扫描文中引用的参考文献，生成参考文献入口文件(aux)文件。
	\item 调用bibtex，以参考文献格式和入口文件为输入，生成格式化以后的参考文献条目文件(bib)。
	\item 再次调用xelatex编译模板，将格式化以后的参考文献条目插入正文。
\end{enumerate}

参考文献数据库(thesis.bib)的条目，可以从Google Scholar搜索引擎\footnote{\url{https://scholar.google.com}}、CiteSeerX搜索引擎\footnote{\url{http://citeseerx.ist.psu.edu}}中查找，文献管理软件Papers\footnote{\url{http://papersapp.com}}、Mendeley\footnote{\url{http://www.mendeley.com}}、JabRef\footnote{\url{http://jabref.sourceforge.net}}也能够输出条目信息。

下面是在Google Scholar上搜索到的一条文献信息，格式是纯文本：

\begin{lstlisting}[caption={从Google Scholar找到的参考文献条目}, label=googlescholar, escapeinside="", numbers=none]
    @phdthesis{"白2008信用风险传染模型和信用衍生品的定价",
      title={"信用风险传染模型和信用衍生品的定价"},
      author={"白云芬"},
      year={2008},
      school={"上海交通大学"}
    }
\end{lstlisting}

推荐修改后在bib文件中的内容为：

\begin{lstlisting}[caption={修改后的参考文献条目}, label=itemok, escapeinside="", numbers=none]
  @phdthesis{bai2008,
    title={"信用风险传染模型和信用衍生品的定价"},
    author={"白云芬"},
    date={2008},
    address={"上海"},
    school={"上海交通大学"}
  }
\end{lstlisting}

按照教务处的要求，参考文献外观应符合国标GBT7714的要求\footnote{\url{http://www.cces.net.cn/guild/sites/tmxb/Files/19798_2.pdf}}。
在模板中，表现形式的控制逻辑通过biblatex-gb7714-2015包实现\footnote{\url{https://www.ctan.org/pkg/biblatex-gb7714-2015}}，基于{Bib\LaTeX}管理文献。在目前的多数TeX发行版中，可能都没有默认包含biblatex-gb7714-2015，需要手动安装。

正文中引用参考文献时，用\verb+\cite{key1,key2,key3...}+可以产生“上标引用的参考文献”，
如\cite{Meta_CN,chen2007act,DPMG}。
使用\verb+\parencite{key1,key2,key3...}+则可以产生水平引用的参考文献，例如\parencite{JohnD,zhubajie,IEEE-1363}。
请看下面的例子，将会穿插使用水平的和上标的参考文献：关于书的\parencite{Meta_CN,JohnD,IEEE-1363}，关于期刊的\cite{chen2007act,chen2007ewi}，
会议论文\parencite{DPMG,kocher99,cnproceed}，
硕士学位论文\parencite{zhubajie,metamori2004}，博士学位论文\cite{shaheshang,FistSystem01,bai2008}，标准文件\parencite{IEEE-1363}，技术报告\cite{NPB2}，电子文献\parencite{xiaoyu2001, CHRISTINE1998}，用户手册\parencite{RManual}。

总结一些注意事项：
\begin{itemize}
\item 参考文献只有在正文中被引用了，才会在最后的参考文献列表中出现；
\item 参考文献“数据库文件”bib是纯文本文件，请使用UTF-8编码，不要使用GBK编码；
\item 参考文献条目中默认通过date域输入时间。兼容使用year域时会产生编译warning，可忽略。
\end{itemize}

\section{用listings插入源代码}

原先ctexbook文档类和listings宏包配合使用时，代码在换页时会出现莫名其妙的错误，后来经高人指点，顺利解决了。
感兴趣的话，可以看看\href{http://bbs.ctex.org/viewthread.php?tid=53451}{这里}。
这里给使用listings宏包插入源代码的例子，这里是一段C代码。
另外，listings宏包真可谓博大精深，可以实现各种复杂、漂亮的效果，想要进一步学习的同学，可以参考
\href{http://mirror.ctan.org/macros/latex/contrib/listings/listings.pdf}{listings宏包手册}。

\begin{lstlisting}[language={C}, caption={一段C源代码}]
#include <stdio.h>
#include <unistd.h>
#include <sys/types.h>
#include <sys/wait.h>

int main() {
  pid_t pid;

  switch ((pid = fork())) {
  case -1:
    printf("fork failed\n");
    break;
  case 0:
    /* child calls exec */
    execl("/bin/ls", "ls", "-l", (char*)0);
    printf("execl failed\n");
    break;
  default:
    /* parent uses wait to suspend execution until child finishes */
    wait((int*)0);
    printf("is completed\n");
    break;
  }

  return 0;
}
\end{lstlisting}

\section{用algorithm和algorithmicx宏包插入算法描述}

algorithmicx 比 algorithmic 增加了一些命令。
示例如算法\ref{algo:sum_100}和算法\ref{algo:merge_sort}，
后者的代码来自\href{http://hustsxh.is-programmer.com/posts/38801.html}{xhSong的博客}。
algorithmicx的详细使用方法见\href{http://mirror.hust.edu.cn/CTAN/macros/latex/contrib/algorithmicx/algorithmicx.pdf}{官方README}。
使用算法宏包时，算法出现的位置很多时候不按照tex文件里的书写顺序,
需要强制定位时可以使用\verb+\begin{algorithm}[H]+
\footnote{http://tex.stackexchange.com/questions/165021/fixing-the-location-of-the-appearance-in-algorithmicx-environment}

这是写在算法\ref{algo:sum_100}前面的一段话，在生成的文件里它会出现在算法\ref{algo:sum_100}前面。

\begin{algorithm}
% \begin{algorithm}[H] % 强制定位
\caption{求100以内的整数和}
\label{algo:sum_100}
\begin{algorithmic}[1] %每行显示行号
\Ensure 100以内的整数和 % 输出
\State $sum \gets 0$
\For{$i = 0 \to 100$}
    \State $sum \gets sum + i$
  \EndFor
\end{algorithmic}
\end{algorithm}

这是写在两个算法中间的一段话，当算法\ref{algo:sum_100}不使用\verb+\begin{algorithm}[H]+时它也会出现在算法\ref{algo:sum_100}前面。

对于很长的算法，单一的算法块\verb+\begin{algorithm}...\end{algorithm}+是不能自动跨页的
\footnote{http://tex.stackexchange.com/questions/70733/latex-algorithm-not-display-under-correct-section}，
会出现的情况有：

\begin{itemize}
  \item 该页放不下当前的算法，留下大片空白，算法在下一页显示
  \item 单一页面放不下当前的算法，显示时超过页码的位置直到超出整个页面范围
\end{itemize}

解决方法有：

\begin{itemize}
  \item (推荐)使用\verb+algstore{algname}+和\verb+algrestore{algname}+来讲算法分为两个部分\footnote{http://tex.stackexchange.com/questions/29816/algorithm-over-2-pages}，如算法\ref{algo:merge_sort}。
  \item 人工拆分算法为多个小的部分。
\end{itemize}

\begin{algorithm}
% \begin{algorithm}[H] % 强制定位
\caption{用归并排序求逆序数}
\label{algo:merge_sort}
\begin{algorithmic}[1] %每行显示行号
\Require $Array$数组，$n$数组大小 % 输入
\Ensure 逆序数 % 输出
\Function {MergerSort}{$Array, left, right$}
  \State $result \gets 0$
  \If {$left < right$}
    \State $middle \gets (left + right) / 2$
    \State $result \gets result +$ \Call{MergerSort}{$Array, left, middle$}
    \State $result \gets result +$ \Call{MergerSort}{$Array, middle, right$}
    \State $result \gets result +$ \Call{Merger}{$Array,left,middle,right$}
  \EndIf
  \State \Return{$result$}
\EndFunction
\State %空一行
\Function{Merger}{$Array, left, middle, right$}
  \State $i\gets left$
  \State $j\gets middle$
  \State $k\gets 0$
  \State $result \gets 0$
  \While{$i<middle$ \textbf{and} $j<right$}
    \If{$Array[i]<Array[j]$}
      \State $B[k++]\gets Array[i++]$
    \Else
      \State $B[k++] \gets Array[j++]$
      \State $result \gets result + (middle - i)$
    \EndIf
  \EndWhile
  \algstore{MergeSort}
\end{algorithmic}
\end{algorithm}

\begin{algorithm}
\begin{algorithmic}[1]
  \algrestore{MergeSort}
  \While{$i<middle$}
    \State $B[k++] \gets Array[i++]$
  \EndWhile
  \While{$j<right$}
    \State $B[k++] \gets Array[j++]$
  \EndWhile
  \For{$i = 0 \to k-1$}
    \State $Array[left + i] \gets B[i]$
  \EndFor
  \State \Return{$result$}
\EndFunction
\end{algorithmic}
\end{algorithm}

这是写在算法\ref{algo:merge_sort}后面的一段话，
但是当算法\ref{algo:merge_sort}不使用\verb+\begin{algorithm}[H]+时它会出现在算法\ref{algo:merge_sort}
甚至算法\ref{algo:sum_100}前面。

对于算法的索引要注意\verb+\caption+和\verb+\label+的位置,
必须是先\verb+\caption+再\verb+\label+\footnote{http://tex.stackexchange.com/questions/65993/algorithm-numbering}，
否则会出现\verb+\ref{algo:sum_100}+生成的编号跟对应算法上显示不一致的问题。

根据Werner的回答\footnote{http://tex.stackexchange.com/questions/53357/switch-cases-in-algorithmic}
增加了\verb+Switch+和\verb+Case+的支持，见算法\ref{algo:switch_example}。

\begin{algorithm}
\caption{Switch示例}
\label{algo:switch_example}
\begin{algorithmic}[1]
  \Switch{$s$}
    \Case{$a$}
      \Assert{0}
    \EndCase
    \Case{$b$}
      \Assert{1}
    \EndCase
    \Default
      \Assert{2}
    \EndDefault
  \EndSwitch
\end{algorithmic}
\end{algorithm}

\include{tex/faq}
\include{tex/summary}

\appendix % 使用英文字母对附录编号

% 附录内容，本科学位论文可以用翻译的文献替代。
\include{tex/app_setup}
\include{tex/app_eq}
\include{tex/app_cjk}
\include{tex/app_log}

\backmatter % 文后无编号部分

% 参考资料
\printbibliography[heading=bibintoc]

% 致谢、发表论文、申请专利、参与项目、简历
% 用于盲审的论文需隐去致谢、发表论文、申请专利、参与的项目
\makeatletter

\ifzstu@coursepaper
\else

  % "研究生学位论文送盲审印刷格式的统一要求"
  % http://www.gs.zstu.edu.cn/inform/3/2015/20151120_123928_738.htm

  % 盲审删去删去致谢页
  \ifzstu@review\relax\else
    \include{tex/ack}         % 致谢
  \fi

  \ifzstu@bachelor
    % 学士学位论文要求在最后有一个英文大摘要，单独编页码
    \include{tex/end_english_abstract}
  \else
    % 盲审论文中，发表学术论文及参与科研情况等仅以第几作者注明即可，不要出现作者或他人姓名
    \ifzstu@review\relax
      \include{tex/pubreview}
      \include{tex/projectsreview}
    \else
      \include{tex/pub}       % 发表论文
      \include{tex/projects}  % 参与的项目
      % \include{tex/patents}   % 申请专利
      \include{tex/resume}    % 个人简历
    \fi
  \fi
\fi

\makeatother

\end{document}
